\section{Experiments}\label{sec:exp}

We ask three questions. (Q1) Is the stack correct and stable in lifelong
operation---no conflicts, no depletions, no deadlocks---across layouts, fleet
sizes and load regimes? (Q2) Does coupling allocation to the planner through
path-aware costs improve throughput or service time over static costs?
(Q3) Which of the remaining design choices (planner backend, window,
replanning period, endpoint capacity, battery term) matter, and what do they
cost in planner time?

\subsection{Experimental Setup}\label{subsec:setup}

\noindent\textbf{Layouts.}
\emph{warehouse} ($52\times19$, 715 free cells) has four blocks of five
10-cell shelf lines with 2-wide aisles, 100 pickup points, 9 docks, 4 charger
stations of 2 slots and 48 home bays. \emph{console} ($36\times25$, 658 free
cells) is generated from the zone/line/dock counts that appear in the
operator console's wireframe labels (zones A--D with 3, 3, 2 and 4 shelf
lines; two named shipping docks, modelled as bays of three delivery cells), 84
pickups, 6 docks, 3 stations, 32 bays. \emph{small} ($16\times10$, 104 free
cells) is an open floor with 4 pickups, 2 docks, 2 slots and 8 bays used for
the planner comparison. \Cref{tab:layouts} in the appendix lists all counts.

\noindent\textbf{Tasks and load.}
Each run delivers 200 tasks (100 on \emph{small}) whose pickup and dock are
sampled uniformly. Two load regimes are used: a Poisson stream of
$\lambda=0.15$ tasks per tick, which is below the capacity of fleets of 16 or
more robots and saturates a fleet of 8, and \emph{batch}, which releases all
tasks at $t=0$ and measures pure throughput. Loading and unloading dwell is
$\delta=2$ ticks. Initial charge is uniform in $[0.5,1.0]$, except in the
battery stress where it is uniform in $[0.2,0.4]$ so that every robot has to
charge during the run. All parameters are listed in \cref{tab:params}.

\noindent\textbf{Methods.}
Four allocation cost models share the same planner, charging policy and
endpoint capacities: \emph{Greedy} (cheapest pair first, static distance),
\emph{Hungarian} (optimal assignment, static distance), \emph{Proxy}
(Hungarian on distance plus reservation-table congestion, \cref{eq:cost}) and
\emph{\sname} (Hungarian on windowed A* arrival times). The planner backend
is prioritized planning (PP) with the hold cascade unless stated; windowed
CBS is compared on \emph{small}. Defaults are $w=20$, $h=5$, $K=6$,
$\kappa_{\mathrm{d}}=3$, $\kappa_{\mathrm{p}}=1$, $w_{\mathrm{b}}=1$.

\noindent\textbf{Metrics.}
Throughput is delivered tasks per 100 ticks; service time is delivery minus
arrival; wait time is first assignment minus arrival; wait actions count
ticks a robot with a goal did not move; holds count planner fall-backs; planner
ms/tick is wall-clock time of allocation plus planning divided by simulated
ticks; conflicts are counted on the executed step (vertex and swap) from the
robots' cells, independently of the reservation table; a run is flagged
deadlocked when no task is picked up or delivered for 1000 consecutive ticks
while tasks are outstanding, and a run that neither finishes nor stalls stops
at 20{,}000 ticks with its unfinished tasks counted. Every configuration is
run with seeds 1--5 and reported as mean\stdv{std}. Runs with the same seed
see the same task stream and initial charge, so we compare variants by
seed-paired relative differences with two-sided 95\% Student-$t$ intervals
(4 degrees of freedom; \cref{tab:paired} in the appendix).

\noindent\textbf{Hardware and protocol.}
All runs are single-threaded Node.js~22 on one core of an AMD Ryzen~5
5500GT; the full benchmark (420 runs) takes 104\,s of wall-clock time. Runs are
deterministic given the seed. The benchmark driver, the per-run JSON records
and the summary CSVs are committed with the code.

\subsection{Main Results}\label{subsec:main}
\begin{table}[t]
\centering
\caption{\textbf{Allocation methods across layouts, fleet sizes and load regimes.} Throughput (delivered tasks per 100 ticks) and mean service time (ticks from arrival to delivery), mean\stdv{std} over 5 seeds, 200 tasks per run, prioritized planning with $w{=}20$, $h{=}5$. Load $\lambda{=}0.15$ is a Poisson stream of 0.15 tasks/tick; \emph{batch} releases all 200 tasks at $t{=}0$. Best mean per row in bold; seed-paired differences with 95\% intervals are in \cref{tab:paired}.}
\label{tab:main}
\small
\setlength{\tabcolsep}{4pt}
\begin{tabular}{llrcccccccc}
\toprule
& & & \multicolumn{4}{c}{\textbf{Throughput} (tasks / 100 ticks) $\uparrow$} & \multicolumn{4}{c}{\textbf{Service time} (ticks) $\downarrow$}\\
\cmidrule(lr){4-7}\cmidrule(lr){8-11}
\textbf{Layout} & \textbf{Load} & $|\mathcal{R}|$ & Greedy & Hungarian & Proxy & \sname & Greedy & Hungarian & Proxy & \sname\\
\midrule
warehouse & $\lambda{=}0.15$ & 8 & 9.3\stdv{0.2} & \textbf{9.3}\stdv{0.3} & 9.3\stdv{0.2} & 9.3\stdv{0.3} & \textbf{312}\stdv{64} & 313\stdv{55} & 315\stdv{56} & 315\stdv{57}\\
warehouse & $\lambda{=}0.15$ & 16 & 13.6\stdv{0.9} & 13.6\stdv{1.0} & 13.5\stdv{1.0} & \textbf{13.6}\stdv{0.9} & 71\stdv{3} & 70\stdv{5} & 74\stdv{8} & \textbf{70}\stdv{6}\\
warehouse & $\lambda{=}0.15$ & 32 & 13.9\stdv{0.9} & 13.9\stdv{0.9} & 13.9\stdv{0.9} & \textbf{13.9}\stdv{0.9} & 53\stdv{1} & \textbf{53}\stdv{1} & 53\stdv{1} & 53\stdv{1}\\
warehouse & batch & 8 & 9.5\stdv{0.3} & 9.5\stdv{0.4} & \textbf{9.5}\stdv{0.3} & 9.5\stdv{0.3} & 963\stdv{25} & 960\stdv{30} & 954\stdv{30} & \textbf{952}\stdv{31}\\
warehouse & batch & 16 & 17.6\stdv{0.7} & 17.5\stdv{0.8} & 17.6\stdv{1.0} & \textbf{17.7}\stdv{0.6} & 496\stdv{23} & 491\stdv{20} & 492\stdv{20} & \textbf{491}\stdv{21}\\
warehouse & batch & 32 & 29.7\stdv{2.3} & 30.6\stdv{2.5} & \textbf{30.6}\stdv{2.2} & 30.3\stdv{2.2} & 287\stdv{12} & 283\stdv{14} & \textbf{283}\stdv{12} & 287\stdv{12}\\
\midrule
console & $\lambda{=}0.15$ & 8 & 11.6\stdv{0.3} & \textbf{11.6}\stdv{0.6} & 11.4\stdv{0.5} & 11.6\stdv{0.5} & 149\stdv{33} & \textbf{149}\stdv{34} & 152\stdv{34} & 150\stdv{33}\\
console & $\lambda{=}0.15$ & 16 & \textbf{13.8}\stdv{0.9} & 13.8\stdv{0.8} & 13.8\stdv{0.9} & 13.8\stdv{0.9} & \textbf{58}\stdv{3} & 59\stdv{3} & 58\stdv{3} & 59\stdv{4}\\
console & $\lambda{=}0.15$ & 32 & 14.0\stdv{0.8} & \textbf{14.0}\stdv{0.8} & 13.9\stdv{0.8} & 14.0\stdv{0.8} & 47\stdv{2} & 47\stdv{2} & \textbf{46}\stdv{3} & 47\stdv{3}\\
console & batch & 8 & 12.2\stdv{0.3} & 12.1\stdv{0.4} & 12.2\stdv{0.5} & \textbf{12.4}\stdv{0.4} & 756\stdv{24} & 759\stdv{21} & 760\stdv{30} & \textbf{752}\stdv{26}\\
console & batch & 16 & 20.0\stdv{0.6} & \textbf{20.1}\stdv{0.8} & 20.0\stdv{0.7} & 19.9\stdv{0.8} & 413\stdv{13} & \textbf{407}\stdv{8} & 411\stdv{10} & 409\stdv{7}\\
console & batch & 32 & 29.2\stdv{2.6} & \textbf{29.8}\stdv{2.2} & 29.3\stdv{2.3} & 29.1\stdv{2.1} & 301\stdv{2} & \textbf{300}\stdv{3} & 301\stdv{8} & 301\stdv{8}\\
\bottomrule
\end{tabular}
\end{table}


\noindent\textbf{Correctness and stability (Q1).}
Over the 420 runs of all sweeps---8{,}833{,}050 robot-ticks, 5{,}317{,}603
move actions and 81{,}000 delivered tasks---the executed steps contain zero
vertex or swap conflicts, as \cref{cor:exec} requires, no robot ever
depletes, and every run delivers all of its tasks without tripping the stall
detector, on every layout, fleet size from 4 to 48, load regime, allocation
method and planner backend. The last two are empirical findings on five
seeds per configuration, not guarantees. Hold fall-backs are rare (at most
5.2 per run on average in the main sweep, reached with 32 robots under batch
load). The battery stress (\cref{tab:battery})
forces 22.6--23.2 charging sessions per run for 16 robots and still records
zero depletions, with a final mean charge of 0.34--0.38.

\noindent\textbf{Allocation coupling (Q2).}
\Cref{tab:main} is the controlled comparison. In none of the twelve
layout/fleet/load cells does the choice of cost model change throughput by
more than 2.9\% (warehouse, 32 robots, batch: 29.7 greedy vs.\ 30.6
Hungarian and Proxy vs.\ 30.3 \sname), and in every cell the four means lie
within one standard deviation of each other; the same holds for service
time. Paired by seed (\cref{tab:paired}), 35 of the 36 throughput intervals
against Hungarian contain zero; the exception is \sname\ on the console with
32 robots under batch load, $2.3\pm2.0\%$ \emph{below} Hungarian, fewer
exclusions than the 1.8 expected by chance among 36 intervals. Every interval
lies within $\pm3.7\%$ under the Poisson load and within $\pm7.1\%$ under
batch load, so the study excludes large effects of the cost model, not effects
of a few percent. This includes the cells where coupling could plausibly help: the
saturated batch regime with 32 robots, where 233--297 wait actions per run
show that robots do interact, and the light regime with 16--32 robots, where
robots are idle and the pairing decision is free to matter. Greedy
assignment is as good as optimal assignment throughout. The reason is
structural rather than statistical, and \cref{sec:disc} discusses it:
the quantity the coupling improves, the empty trip to the pickup, is a small
and well-predicted part of a task, and the congestion that exists arises at
endpoints rather than along paths.

\noindent\textbf{Planner backend (Q3, part 1).}
\Cref{tab:planner} and \cref{fig:planner} compare PP and windowed CBS on
\emph{small}, where CBS is affordable. CBS shortens the makespan by 2.8\%,
8.3\% and 8.3\% for 4, 6 and 8 robots (seed-paired: $-2.8\pm2.1$,
$-8.2\pm3.6$ and $-8.2\pm4.3\%$), cuts wait actions by 19--31\% and
service time by 6--10\%, at 0.10--0.32 planner ms per tick; it exceeded its
node budget in 0.6 epochs per run only at 6 robots, falling back to PP for
that group. The gain is exactly the suboptimality of fixed priorities
inside the window; it does not carry to the larger layouts, where a group of
32 busy robots is beyond a 2000-node budget in most epochs.
\begin{table}[t]
\centering
\caption{\textbf{Planner backend on the small layout} (batch load, 100 tasks, Hungarian allocation, 5 seeds). Windowed CBS resolves conflicts inside the window optimally and falls back to prioritized planning (PP) when its node budget (2000) is exceeded.}
\label{tab:planner}
\small
\begin{tabular}{rlccccc}
\toprule
$|\mathcal{R}|$ & \textbf{Planner} & \textbf{Makespan} $\downarrow$ & \textbf{Service time} $\downarrow$ & \textbf{Wait actions} $\downarrow$ & \textbf{Planner ms/tick} & \textbf{CBS fallbacks}\\
\midrule
4 & PP & 789\stdv{39} & 318.1\stdv{10.3} & 60\stdv{13} & 0.08\stdv{0.01} & --\\
 & Windowed CBS & \textbf{767}\stdv{39} & \textbf{300}\stdv{11} & \textbf{48}\stdv{6} & 0.09\stdv{0.01} & 0.0\stdv{0.0}\\
6 & PP & 562\stdv{37} & 262.0\stdv{8.0} & 135\stdv{17} & 0.19\stdv{0.02} & --\\
 & Windowed CBS & \textbf{515}\stdv{27} & \textbf{239}\stdv{10} & \textbf{102}\stdv{10} & 0.31\stdv{0.07} & 0.6\stdv{0.5}\\
8 & PP & 513\stdv{31} & 242.1\stdv{9.4} & 214\stdv{35} & 0.27\stdv{0.02} & --\\
 & Windowed CBS & \textbf{471}\stdv{28} & \textbf{218}\stdv{13} & \textbf{148}\stdv{16} & 0.34\stdv{0.04} & 0.0\stdv{0.0}\\
\bottomrule
\end{tabular}
\end{table}

\begin{figure}[t]
\centering
\begin{tikzpicture}
\begin{groupplot}[
  group style={group size=3 by 1, horizontal sep=14mm},
  width=0.34\linewidth, height=42mm,
  tick label style={font=\scriptsize}, label style={font=\scriptsize},
  legend style={font=\scriptsize, draw=none, fill=none},
  error bars/y dir=both, error bars/y explicit,
  grid=major, grid style={gray!25},
]
\nextgroupplot[xlabel={fleet size $|\mathcal{R}|$ (small layout)}, ylabel={makespan (ticks)}, xtick={4,6,8}, legend pos=north east]
\addplot[blue, mark=*, mark size=1.4pt] table[col sep=comma, x=fleetSize, y=makespan_mean, y error=makespan_std] {figures/data/planner_pp.csv};
\addlegendentry{PP}
\addplot[red, mark=square*, mark size=1.4pt] table[col sep=comma, x=fleetSize, y=makespan_mean, y error=makespan_std] {figures/data/planner_cbs.csv};
\addlegendentry{windowed CBS}
\nextgroupplot[xlabel={fleet size $|\mathcal{R}|$ (small layout)}, ylabel={wait actions per run}, xtick={4,6,8}]
\addplot[blue, mark=*, mark size=1.4pt] table[col sep=comma, x=fleetSize, y=waitActions_mean, y error=waitActions_std] {figures/data/planner_pp.csv};
\addplot[red, mark=square*, mark size=1.4pt] table[col sep=comma, x=fleetSize, y=waitActions_mean, y error=waitActions_std] {figures/data/planner_cbs.csv};
\nextgroupplot[xlabel={window $w$ (warehouse, 32 robots)}, ylabel={throughput / planner ms per tick}, xtick={10,20,30}, legend pos=north west]
\addplot[black, mark=*, mark size=1.4pt] table[col sep=comma, x=window, y=throughput_mean, y error=throughput_std] {figures/data/window.csv};
\addlegendentry{throughput}
\addplot[orange, mark=triangle*, mark size=1.6pt] table[col sep=comma, x=window, y expr=10*\thisrow{plannerMsPerTick_mean}, y error expr=10*\thisrow{plannerMsPerTick_std}] {figures/data/window.csv};
\addlegendentry{$10\times$ planner ms/tick}
\end{groupplot}
\end{tikzpicture}
\caption{\textbf{Planner backend and window.} Left, centre: on the small layout, windowed CBS shortens the makespan and removes wait actions relative to prioritized planning (PP) at every fleet size (\cref{tab:planner}). Right: on the warehouse with 32 robots, a shorter window raises throughput and cuts planner time (\cref{tab:ablation}); the planner cost is scaled by 10 to share the axis.}
\label{fig:planner}
\end{figure}


\noindent\textbf{Scaling (Q3, part 2).}
\Cref{fig:scaling} (data in \cref{tab:scaling}) sweeps the warehouse fleet
from 4 to 48 robots under batch load. Throughput rises from 4.9 to 33.7 tasks
per 100 ticks and flattens beyond 40 robots, where utilisation has fallen to
0.41 because $9$ docks $\times\,\kappa_{\mathrm{d}}=3$ bound the number of
tasks in flight; service time falls from 1909 to 252 ticks. Planner cost
grows roughly linearly, from 0.02 to 0.77 ms per simulated tick for Hungarian
and to 1.00 for \sname; the path-aware pricing costs 35\% more A* calls at
48 robots (16{,}229 vs.\ 12{,}021 per run) and buys nothing measurable.
\begin{figure}[t]
\centering
\begin{tikzpicture}
\begin{groupplot}[
  group style={group size=3 by 1, horizontal sep=14mm},
  width=0.36\linewidth, height=42mm,
  xlabel={fleet size $|\mathcal{R}|$}, xtick={4,8,16,24,32,40,48},
  tick label style={font=\scriptsize}, label style={font=\scriptsize},
  legend style={font=\scriptsize, draw=none, fill=none},
  error bars/y dir=both, error bars/y explicit,
  grid=major, grid style={gray!25},
]
\nextgroupplot[ylabel={throughput (tasks / 100 ticks)}, legend to name=scalinglegend, legend columns=2]
\addplot[blue, mark=*, mark size=1.4pt] table[col sep=comma, x=fleetSize, y=throughput_mean, y error=throughput_std] {figures/data/scaling_hungarian.csv};
\addlegendentry{Hungarian (static cost)}
\addplot[red, mark=square*, mark size=1.4pt] table[col sep=comma, x=fleetSize, y=throughput_mean, y error=throughput_std] {figures/data/scaling_pact.csv};
\addlegendentry{\sname\ (path-aware cost)}
\nextgroupplot[ylabel={planner ms per tick}]
\addplot[blue, mark=*, mark size=1.4pt] table[col sep=comma, x=fleetSize, y=plannerMsPerTick_mean, y error=plannerMsPerTick_std] {figures/data/scaling_hungarian.csv};
\addplot[red, mark=square*, mark size=1.4pt] table[col sep=comma, x=fleetSize, y=plannerMsPerTick_mean, y error=plannerMsPerTick_std] {figures/data/scaling_pact.csv};
\nextgroupplot[ylabel={A* calls per run}]
\addplot[blue, mark=*, mark size=1.4pt] table[col sep=comma, x=fleetSize, y=astarCalls_mean, y error=astarCalls_std] {figures/data/scaling_hungarian.csv};
\addplot[red, mark=square*, mark size=1.4pt] table[col sep=comma, x=fleetSize, y=astarCalls_mean, y error=astarCalls_std] {figures/data/scaling_pact.csv};
\end{groupplot}
\node at ($(group c2r1.south)+(0,-11mm)$) {\pgfplotslegendfromname{scalinglegend}};
\end{tikzpicture}
\caption{\textbf{Fleet-size scaling} on the warehouse layout under batch load (200 tasks, mean\stdv{std} over 5 seeds; data in \cref{tab:scaling}). Throughput saturates once dock capacity ($9$ docks $\times\,\kappa_{\mathrm{d}}{=}3$) binds; planner cost grows roughly linearly in the fleet and stays below one millisecond per simulated tick; path-aware pricing adds the extra A* calls of the third panel.}
\label{fig:scaling}
\end{figure}


\subsection{Ablation Studies}\label{subsec:ablation}
\begin{table}[t]
\centering
\caption{\textbf{Ablations} on the warehouse layout with 32 robots under batch load (200 tasks, 5 seeds). Each row changes one setting of the default configuration.}
\label{tab:ablation}
\scriptsize
\begin{tabular}{lccccc}
\toprule
\textbf{Variant} & \textbf{Throughput} $\uparrow$ & \textbf{Service time} $\downarrow$ & \textbf{Wait actions} & \textbf{Holds} & \textbf{Planner ms/tick}\\
\midrule
\sname\ (default: $w{=}20$, $h{=}5$, $K{=}6$, $\kappa_{\mathrm{d}}{=}3$, $w_{\mathrm{b}}{=}1$) & 30.3\stdv{2.2} & 287\stdv{12} & 274\stdv{60} & 4.0\stdv{2.8} & 0.61\stdv{0.05}\\
\quad no battery term ($w_{\mathrm{b}}{=}0$) & 30.3\stdv{2.3} & 285\stdv{13} & 253\stdv{21} & 2.8\stdv{1.8} & 0.60\stdv{0.07}\\
\quad Hungarian, no battery term & 30.1\stdv{1.8} & 283\stdv{9} & 253\stdv{32} & 3.0\stdv{2.3} & 0.57\stdv{0.07}\\
\quad candidates $K{=}3$ & 30.2\stdv{2.3} & 285\stdv{13} & 262\stdv{28} & 1.8\stdv{1.8} & 0.60\stdv{0.05}\\
\quad candidates $K{=}12$ & 30.4\stdv{2.0} & 284\stdv{14} & 228\stdv{33} & 2.4\stdv{1.1} & 0.62\stdv{0.07}\\
\quad window $w{=}10$ & 32.4\stdv{2.1} & 261\stdv{11} & 150\stdv{7} & 1.2\stdv{1.3} & 0.19\stdv{0.00}\\
\quad window $w{=}30$ & 28.2\stdv{2.6} & 307\stdv{15} & 357\stdv{72} & 2.8\stdv{1.3} & 1.96\stdv{0.16}\\
\quad replan every tick ($h{=}1$) & 30.3\stdv{2.3} & 287\stdv{12} & 257\stdv{34} & 4.0\stdv{0.7} & 1.35\stdv{0.15}\\
\quad no event-triggered epochs & 29.9\stdv{2.3} & 288\stdv{12} & 329\stdv{45} & 0.8\stdv{0.8} & 0.30\stdv{0.01}\\
\quad Hungarian, no event-triggered epochs & 28.9\stdv{1.7} & 290\stdv{14} & 359\stdv{35} & 1.4\stdv{0.5} & 0.26\stdv{0.02}\\
\quad dock capacity $\kappa_{\mathrm{d}}{=}2$ & 24.0\stdv{1.5} & 344\stdv{16} & 125\stdv{14} & 1.4\stdv{1.1} & 0.39\stdv{0.02}\\
\quad Hungarian, dock capacity $\kappa_{\mathrm{d}}{=}2$ & 23.7\stdv{1.8} & 347\stdv{16} & 127\stdv{12} & 1.0\stdv{0.7} & 0.30\stdv{0.02}\\
\bottomrule
\end{tabular}
\end{table}


\Cref{tab:ablation} varies one setting at a time on the warehouse with 32
robots under batch load.

\noindent\textbf{Endpoint capacity is the dominant lever.}
Lowering the per-dock in-flight capacity from 3 to 2 drops throughput from
30.3 to 24.0 (\sname) and from 30.6 to 23.7 (Hungarian, \cref{tab:main}),
i.e.\ capacity 3 is worth $+26$--$29\%$ (seed-paired, capacity 2 costs
$20.7\pm3.3\%$ and $22.5\pm2.2\%$), and utilisation falls from 0.64 to
0.42 because tasks wait in the pool for a free dock. During development,
capacity 4 produced dock-front gridlocks in 2 of 30 runs under an earlier
priority scheme; these runs are not part of the committed benchmark.
Capacity 3 completed every committed run and is the default.

\noindent\textbf{Window and period trade throughput for planner time.}
A window of 10 ticks raises throughput to 32.4 ($+7\%$, seed-paired
$+6.9\pm1.6\%$) and halves wait actions (150 vs.\ 274) while cutting
planner time from 0.63 to 0.18 ms/tick; a window of 30 lowers throughput to
28.2 ($-6.9\pm3.9\%$) at 1.98 ms/tick. Longer
windows commit robots to longer reserved waits behind others and make
fall-backs more likely, the same effect that \citet{li2021rhcr} report for
RHCR. Replanning every tick ($h=1$) leaves throughput unchanged (30.3) at
$2.2\times$ the planner time. Disabling event-triggered epochs, so that idle
robots wait for the next periodic epoch, costs 1--5\% throughput and adds
55--125 wait actions; only the Hungarian loss is resolved by the paired
interval ($-5.2\pm3.3\%$, against $-1.2\pm5.2\%$ for \sname).

\noindent\textbf{Battery term and candidate pruning are neutral.}
Removing the battery factor ($w_{\mathrm{b}}=0$) changes throughput by less
than 0.5 for both cost models ($+0.01$ for \sname, $-0.43$ for Hungarian),
also in the battery stress (\cref{tab:battery}); with $\theta_{\mathrm{lo}}=0.2$ and the feasibility
check of \cref{eq:feasible}, charging decisions are rarely marginal. Pricing
3 or 12 candidates instead of 6 changes throughput by at most 0.2.
